\section{Four rows}\label{sec:four}

On four rows the row bound of Theorem~\ref{thm:rowbound} is $m+4$. It is attained for every $m\ge6$, and the widths $4$ and $5$ fall short by one. Four rows differ from three in the count of optima. On three rows a fixed fraction $4/27$ of the axis-valid candidates survives the tilted rectangles. On four rows the surviving fraction tends to zero exponentially. The optima are still recognized by a finite machine, so their generating function is rational, but the machine remembers six columns, and four would not be enough. The last subsection computes the generating function, and $a_4(m)$ grows like $m^4$ times $3.3520^m$. Throughout, $(m)_k=m(m-1)\cdots(m-k+1)$, and the tilted rectangles of the $4\times m$ grid are the diamonds $D(i,1)$, $D(i,2)$ and the kinds $\kA_i$, $\kB_i$, $\kC_i$, $\kE_i$ of Lemma~\ref{lem:tilted-narrow}.

\subsection{The value for \texorpdfstring{$m\ge6$}{m at least 6}}

\begin{theorem}[four rows]\label{thm:fourvalue}
For every $m\ge6$ we have $f(4,m)=m+4$.
\end{theorem}

The upper bound is Theorem~\ref{thm:rowbound} with $h=4$. For the lower bound we take
\begin{equation}\label{eq:fourS}
S=\{(i,1):2\le i\le m-1\}\cup\{(m-1,2),(m,2),(1,3),(2,3),(1,4),(m,4)\},
\end{equation}
which has $m+4$ points. Figure~\ref{fig:four} shows it for $m=7$. The double columns are column $1$ (rows $3,4$), column $2$ (rows $1,3$), column $m-1$ (rows $1,2$), and column $m$ (rows $2,4$), and every other column is black in row $1$ only. So $G_S$ is the eight-cycle $r_1\,c_2\,r_3\,c_1\,r_4\,c_m\,r_2\,c_{m-1}$ with the columns $c_3,\ldots,c_{m-2}$ hanging off $r_1$.

\begin{figure}[t]
\centering
\begin{tikzpicture}[x=0.520cm,y=0.520cm,baseline=(current bounding box.north)]
  \draw[gridcol] (1,-1) grid (7,-4);
  \draw[cyccol,line width=1.8pt,line join=round,opacity=0.85] (2,-1) -- (6,-1) -- (6,-2) -- (7,-2) -- (7,-4) -- (1,-4) -- (1,-3) -- (2,-3) -- cycle;
  \filldraw[fill=white,draw=keptline,line width=0.6pt] (1,-1) circle (0.17);
  \fill[black] (2,-1) circle (0.17);
  \fill[black] (3,-1) circle (0.17);
  \fill[black] (4,-1) circle (0.17);
  \fill[black] (5,-1) circle (0.17);
  \fill[black] (6,-1) circle (0.17);
  \filldraw[fill=white,draw=keptline,line width=0.6pt] (7,-1) circle (0.17);
  \filldraw[fill=white,draw=keptline,line width=0.6pt] (1,-2) circle (0.17);
  \filldraw[fill=white,draw=keptline,line width=0.6pt] (2,-2) circle (0.17);
  \filldraw[fill=white,draw=keptline,line width=0.6pt] (3,-2) circle (0.17);
  \filldraw[fill=white,draw=keptline,line width=0.6pt] (4,-2) circle (0.17);
  \filldraw[fill=white,draw=keptline,line width=0.6pt] (5,-2) circle (0.17);
  \fill[black] (6,-2) circle (0.17);
  \fill[black] (7,-2) circle (0.17);
  \fill[black] (1,-3) circle (0.17);
  \fill[black] (2,-3) circle (0.17);
  \filldraw[fill=white,draw=keptline,line width=0.6pt] (3,-3) circle (0.17);
  \filldraw[fill=white,draw=keptline,line width=0.6pt] (4,-3) circle (0.17);
  \filldraw[fill=white,draw=keptline,line width=0.6pt] (5,-3) circle (0.17);
  \filldraw[fill=white,draw=keptline,line width=0.6pt] (6,-3) circle (0.17);
  \filldraw[fill=white,draw=keptline,line width=0.6pt] (7,-3) circle (0.17);
  \fill[black] (1,-4) circle (0.17);
  \filldraw[fill=white,draw=keptline,line width=0.6pt] (2,-4) circle (0.17);
  \filldraw[fill=white,draw=keptline,line width=0.6pt] (3,-4) circle (0.17);
  \filldraw[fill=white,draw=keptline,line width=0.6pt] (4,-4) circle (0.17);
  \filldraw[fill=white,draw=keptline,line width=0.6pt] (5,-4) circle (0.17);
  \filldraw[fill=white,draw=keptline,line width=0.6pt] (6,-4) circle (0.17);
  \fill[black] (7,-4) circle (0.17);
  \node[lab] at (1,-0.35) {1};
  \node[lab] at (2,-0.35) {2};
  \node[lab] at (3,-0.35) {3};
  \node[lab] at (4,-0.35) {4};
  \node[lab] at (5,-0.35) {5};
  \node[lab] at (6,-0.35) {6};
  \node[lab] at (7,-0.35) {7};
  \node[lab] at (0.35,-1) {1};
  \node[lab] at (0.35,-2) {2};
  \node[lab] at (0.35,-3) {3};
  \node[lab] at (0.35,-4) {4};
  \node[capt,anchor=north,text width=4.41cm,align=center] at (4,-5.4) {(a) the blackout, $m=7$};
  \draw[draw=cyccol,line width=2pt] (14.155371111771444,-2.2729014785063293) -- (14.155371111771444,-3.7270985214936707);
  \draw[draw=cyccol,line width=2pt] (13.127098521493672,-4.755371111771445) -- (14.155371111771444,-3.7270985214936707);
  \draw[draw=cyccol,line width=2pt] (11.67290147850633,-1.2446288882285552) -- (13.127098521493672,-1.2446288882285552);
  \draw[draw=cyccol,line width=2pt] (14.155371111771444,-2.2729014785063293) -- (13.127098521493672,-1.2446288882285552);
  \draw[draw=pencol,line width=1.3pt,dash pattern=on 3pt off 2pt] (11.67290147850633,-1.2446288882285552) -- (9.372901478506328,-0.3446288882285552);
  \draw[draw=pencol,line width=1.3pt,dash pattern=on 3pt off 2pt] (11.67290147850633,-1.2446288882285552) -- (10.47290147850633,0.055371111771444825);
  \draw[draw=pencol,line width=1.3pt,dash pattern=on 3pt off 2pt] (11.67290147850633,-1.2446288882285552) -- (11.57290147850633,0.2553711117714448);
  \draw[draw=cyccol,line width=2pt] (11.67290147850633,-1.2446288882285552) -- (10.644628888228555,-2.2729014785063297);
  \draw[draw=cyccol,line width=2pt] (10.644628888228556,-3.7270985214936707) -- (10.644628888228555,-2.2729014785063297);
  \draw[draw=cyccol,line width=2pt] (10.644628888228556,-3.7270985214936707) -- (11.67290147850633,-4.755371111771445);
  \draw[draw=cyccol,line width=2pt] (13.127098521493672,-4.755371111771445) -- (11.67290147850633,-4.755371111771445);
  \node[rowv] at (11.67290147850633,-1.2446288882285552) {$r_{1}$};
  \node[colv] at (13.127098521493672,-1.2446288882285552) {$c_{2}$};
  \node[rowv] at (14.155371111771444,-2.2729014785063293) {$r_{3}$};
  \node[colv] at (14.155371111771444,-3.7270985214936707) {$c_{1}$};
  \node[rowv] at (13.127098521493672,-4.755371111771445) {$r_{4}$};
  \node[colv] at (11.67290147850633,-4.755371111771445) {$c_{7}$};
  \node[rowv] at (10.644628888228556,-3.7270985214936707) {$r_{2}$};
  \node[colv] at (10.644628888228555,-2.2729014785063297) {$c_{6}$};
  \node[colv] at (9.372901478506328,-0.3446288882285552) {$c_{3}$};
  \node[colv] at (10.47290147850633,0.055371111771444825) {$c_{4}$};
  \node[colv] at (11.57290147850633,0.2553711117714448) {$c_{5}$};
  \node[capt,anchor=north] at (12.4,-5.4) {(b) $G_S$};
\end{tikzpicture}

\caption{The four-row construction for $m=7$. (a) The blackout. The bold path turns only at black points and alternates between rows and columns. (b) The row-column graph $G_S$. The path of (a) is its eight-cycle, and the columns $3$, $4$, $5$ hang off $r_1$ (dashed).}
\label{fig:four}
\end{figure}

\begin{proof}[Proof of Theorem~\ref{thm:fourvalue}]
\emph{Step 1. Axis-aligned rectangles.} The graph $G_S$ has no cycle of length four or six, so by Lemma~\ref{lem:girth} the axis-aligned rectangles have distinct presentations, and none is hidden.

\emph{Step 2. At most one black corner.} Such a rectangle is safe by Lemma~\ref{lem:threecorners}. In rows $2$ to $4$ the black points are $(m-1,2)$, $(m,2)$, $(1,3)$, $(2,3)$, $(1,4)$, $(m,4)$. The diamond $D(i,2)=\{(i,3),(i+1,2),(i+2,3),(i+1,4)\}$ has $(i,3)$ black only if $i\le2$ and $(i+1,2)$ black only if $i=m-2$, which exclude each other as $m\ge6$, and its other two corners are white. The rectangles $\kC_i=\{(i,2),(i+1,4),(i+2,1),(i+3,3)\}$ and $\kE_i=\{(i,2),(i+2,4),(i+1,1),(i+3,3)\}$ have their corners in rows $2$, $3$, $4$ white, since $i\le m-3$, $2\le i+1\le i+2\le m-2$, and $i+3\ge4$.

\emph{Step 3. The weak tilted rectangles.} For $D(i,1)=\{(i,2),(i+1,1),(i+2,2),(i+1,3)\}$ the corner $(i+1,1)$ is always black, $(i,2)$ never, $(i+2,2)$ exactly when $i\in\{m-3,m-2\}$, and $(i+1,3)$ exactly when $i=1$. For $\kA_i=\{(i,3),(i+1,4),(i+2,1),(i+3,2)\}$ and $\kB_i=\{(i,3),(i+2,4),(i+1,1),(i+3,2)\}$ the corner in row $1$ is always black, the corner in row $4$ never, $(i,3)$ exactly when $i\le2$, and $(i+3,2)$ exactly when $i\in\{m-4,m-3\}$. So the tilted rectangles showing at most two corners are these, with $s=1$ for $\kA_i$ and $s=2$ for $\kB_i$.
\begin{center}\small
\begin{tabular}{@{}lll@{}}
\toprule
rectangle & which $i$ & corners shown\\
\midrule
$D(1,1)$ & & $(1,2),(3,2)$\\
$D(i,1)$ & $i\in\{m-3,m-2\}$ & $(i,2),(i+1,3)$\\
$\kA_i$ or $\kB_i$ & $i\in\{1,2\}$, $i\notin\{m-4,m-3\}$ & $(i+s,4),(i+3,2)$\\
$\kA_i$ or $\kB_i$ & $i\in\{m-4,m-3\}$, $i\notin\{1,2\}$ & $(i,3),(i+s,4)$\\
$\kA_i$ or $\kB_i$ & $i\in\{1,2\}\cap\{m-4,m-3\}$ & $(i+s,4)$\\
\bottomrule
\end{tabular}
\end{center}
The last line occurs only for $m=6$ and $i=2$. Figure~\ref{fig:fourtilted} shows five of these rectangles for $m=7$.

\emph{Step 4. Against each other.} The rows of the two shown points are $2,2$ and $2,3$ and $4,2$ and $3,4$ in the first four lines, so different lines never agree. Within a line the point in row $2$ or $3$ determines $i$, and the point in row $4$ then determines $s$. The two presentations of the last line are $(3,4)$ and $(4,4)$.

\emph{Step 5. Against axis-aligned rectangles.} An axis-aligned rectangle showing $(1,2)$ and $(3,2)$ needs a row $r$ with $(1,r)$ and $(3,r)$ black, but column $3$ is black only in row $1$ and $(1,1)$ is white. In the second line the axis-aligned rectangle through $(i,2)$ and $(i+1,3)$ has the white corner $(i,3)$, as $i\ge m-3\ge3$. In the third line its corner $(i+s,2)$ is white, as $i+s\le4<m-1$. In the fourth line its corner $(i,4)$ is white, as $3\le i\le m-3$. In the last line a rival of $\{(z,4)\}$ needs a column $c\in\{1,6\}$ and a row $r$ with $(z,r)$ and $(c,r)$ black, so $r=1$, but $(1,1)$ and $(6,1)$ are white.

\emph{Step 6. Nothing is hidden.} A hidden axis-aligned rectangle would be a four-cycle of $G_S$, and every tilted rectangle shows a corner by Steps 2 and 3. So $S$ is valid.
\end{proof}

\begin{figure}[t]
\centering
\begin{tikzpicture}[x=0.430cm,y=0.430cm,baseline=(current bounding box.north)]
  \draw[gridcol] (1.0,-1) grid (7.0,-4);
  \filldraw[fill=white,draw=keptline,line width=0.6pt] (1.0,-1) circle (0.17);
  \fill[black] (2.0,-1) circle (0.17);
  \fill[black] (3.0,-1) circle (0.17);
  \fill[black] (4.0,-1) circle (0.17);
  \fill[black] (5.0,-1) circle (0.17);
  \fill[black] (6.0,-1) circle (0.17);
  \filldraw[fill=white,draw=keptline,line width=0.6pt] (7.0,-1) circle (0.17);
  \filldraw[fill=white,draw=keptline,line width=0.6pt] (1.0,-2) circle (0.17);
  \filldraw[fill=white,draw=keptline,line width=0.6pt] (2.0,-2) circle (0.17);
  \filldraw[fill=white,draw=keptline,line width=0.6pt] (3.0,-2) circle (0.17);
  \filldraw[fill=white,draw=keptline,line width=0.6pt] (4.0,-2) circle (0.17);
  \filldraw[fill=white,draw=keptline,line width=0.6pt] (5.0,-2) circle (0.17);
  \fill[black] (6.0,-2) circle (0.17);
  \fill[black] (7.0,-2) circle (0.17);
  \fill[black] (1.0,-3) circle (0.17);
  \fill[black] (2.0,-3) circle (0.17);
  \filldraw[fill=white,draw=keptline,line width=0.6pt] (3.0,-3) circle (0.17);
  \filldraw[fill=white,draw=keptline,line width=0.6pt] (4.0,-3) circle (0.17);
  \filldraw[fill=white,draw=keptline,line width=0.6pt] (5.0,-3) circle (0.17);
  \filldraw[fill=white,draw=keptline,line width=0.6pt] (6.0,-3) circle (0.17);
  \filldraw[fill=white,draw=keptline,line width=0.6pt] (7.0,-3) circle (0.17);
  \fill[black] (1.0,-4) circle (0.17);
  \filldraw[fill=white,draw=keptline,line width=0.6pt] (2.0,-4) circle (0.17);
  \filldraw[fill=white,draw=keptline,line width=0.6pt] (3.0,-4) circle (0.17);
  \filldraw[fill=white,draw=keptline,line width=0.6pt] (4.0,-4) circle (0.17);
  \filldraw[fill=white,draw=keptline,line width=0.6pt] (5.0,-4) circle (0.17);
  \filldraw[fill=white,draw=keptline,line width=0.6pt] (6.0,-4) circle (0.17);
  \fill[black] (7.0,-4) circle (0.17);
  \node[lab] at (1.0,-0.35) {1};
  \node[lab] at (2.0,-0.35) {2};
  \node[lab] at (3.0,-0.35) {3};
  \node[lab] at (4.0,-0.35) {4};
  \node[lab] at (5.0,-0.35) {5};
  \node[lab] at (6.0,-0.35) {6};
  \node[lab] at (7.0,-0.35) {7};
  \node[lab] at (0.35,-1) {1};
  \node[lab] at (0.35,-2) {2};
  \node[lab] at (0.35,-3) {3};
  \node[lab] at (0.35,-4) {4};
  \draw[draw=rectA,line width=1.3pt,line join=round] (2.0,-1) -- (3.0,-2) -- (2.0,-3) -- (1.0,-2) -- cycle;
  \draw[rectA,line width=1.1pt] (1.0,-2) circle (0.3);
  \draw[rectA,line width=1.1pt] (3.0,-2) circle (0.3);
  \node[capt,anchor=north] at (4.0,-4.7) {(a) $D(1,1)$};
  \draw[gridcol] (9.5,-1) grid (15.5,-4);
  \filldraw[fill=white,draw=keptline,line width=0.6pt] (9.5,-1) circle (0.17);
  \fill[black] (10.5,-1) circle (0.17);
  \fill[black] (11.5,-1) circle (0.17);
  \fill[black] (12.5,-1) circle (0.17);
  \fill[black] (13.5,-1) circle (0.17);
  \fill[black] (14.5,-1) circle (0.17);
  \filldraw[fill=white,draw=keptline,line width=0.6pt] (15.5,-1) circle (0.17);
  \filldraw[fill=white,draw=keptline,line width=0.6pt] (9.5,-2) circle (0.17);
  \filldraw[fill=white,draw=keptline,line width=0.6pt] (10.5,-2) circle (0.17);
  \filldraw[fill=white,draw=keptline,line width=0.6pt] (11.5,-2) circle (0.17);
  \filldraw[fill=white,draw=keptline,line width=0.6pt] (12.5,-2) circle (0.17);
  \filldraw[fill=white,draw=keptline,line width=0.6pt] (13.5,-2) circle (0.17);
  \fill[black] (14.5,-2) circle (0.17);
  \fill[black] (15.5,-2) circle (0.17);
  \fill[black] (9.5,-3) circle (0.17);
  \fill[black] (10.5,-3) circle (0.17);
  \filldraw[fill=white,draw=keptline,line width=0.6pt] (11.5,-3) circle (0.17);
  \filldraw[fill=white,draw=keptline,line width=0.6pt] (12.5,-3) circle (0.17);
  \filldraw[fill=white,draw=keptline,line width=0.6pt] (13.5,-3) circle (0.17);
  \filldraw[fill=white,draw=keptline,line width=0.6pt] (14.5,-3) circle (0.17);
  \filldraw[fill=white,draw=keptline,line width=0.6pt] (15.5,-3) circle (0.17);
  \fill[black] (9.5,-4) circle (0.17);
  \filldraw[fill=white,draw=keptline,line width=0.6pt] (10.5,-4) circle (0.17);
  \filldraw[fill=white,draw=keptline,line width=0.6pt] (11.5,-4) circle (0.17);
  \filldraw[fill=white,draw=keptline,line width=0.6pt] (12.5,-4) circle (0.17);
  \filldraw[fill=white,draw=keptline,line width=0.6pt] (13.5,-4) circle (0.17);
  \filldraw[fill=white,draw=keptline,line width=0.6pt] (14.5,-4) circle (0.17);
  \fill[black] (15.5,-4) circle (0.17);
  \draw[draw=rectA,line width=1.3pt,line join=round] (11.5,-1) -- (12.5,-2) -- (10.5,-4) -- (9.5,-3) -- cycle;
  \draw[rectA,line width=1.1pt] (10.5,-4) circle (0.3);
  \draw[rectA,line width=1.1pt] (12.5,-2) circle (0.3);
  \draw[draw=rectB,line width=1.3pt,line join=round,dash pattern=on 3pt off 2pt] (10.5,-1) -- (12.5,-2) -- (11.5,-4) -- (9.5,-3) -- cycle;
  \draw[rectB,line width=1.1pt] (11.5,-4) circle (0.38);
  \draw[rectB,line width=1.1pt] (12.5,-2) circle (0.38);
  \node[capt,anchor=north] at (12.5,-4.7) {(b) $\kA_1$ and $\kB_1$};
  \draw[gridcol] (18.0,-1) grid (24.0,-4);
  \filldraw[fill=white,draw=keptline,line width=0.6pt] (18.0,-1) circle (0.17);
  \fill[black] (19.0,-1) circle (0.17);
  \fill[black] (20.0,-1) circle (0.17);
  \fill[black] (21.0,-1) circle (0.17);
  \fill[black] (22.0,-1) circle (0.17);
  \fill[black] (23.0,-1) circle (0.17);
  \filldraw[fill=white,draw=keptline,line width=0.6pt] (24.0,-1) circle (0.17);
  \filldraw[fill=white,draw=keptline,line width=0.6pt] (18.0,-2) circle (0.17);
  \filldraw[fill=white,draw=keptline,line width=0.6pt] (19.0,-2) circle (0.17);
  \filldraw[fill=white,draw=keptline,line width=0.6pt] (20.0,-2) circle (0.17);
  \filldraw[fill=white,draw=keptline,line width=0.6pt] (21.0,-2) circle (0.17);
  \filldraw[fill=white,draw=keptline,line width=0.6pt] (22.0,-2) circle (0.17);
  \fill[black] (23.0,-2) circle (0.17);
  \fill[black] (24.0,-2) circle (0.17);
  \fill[black] (18.0,-3) circle (0.17);
  \fill[black] (19.0,-3) circle (0.17);
  \filldraw[fill=white,draw=keptline,line width=0.6pt] (20.0,-3) circle (0.17);
  \filldraw[fill=white,draw=keptline,line width=0.6pt] (21.0,-3) circle (0.17);
  \filldraw[fill=white,draw=keptline,line width=0.6pt] (22.0,-3) circle (0.17);
  \filldraw[fill=white,draw=keptline,line width=0.6pt] (23.0,-3) circle (0.17);
  \filldraw[fill=white,draw=keptline,line width=0.6pt] (24.0,-3) circle (0.17);
  \fill[black] (18.0,-4) circle (0.17);
  \filldraw[fill=white,draw=keptline,line width=0.6pt] (19.0,-4) circle (0.17);
  \filldraw[fill=white,draw=keptline,line width=0.6pt] (20.0,-4) circle (0.17);
  \filldraw[fill=white,draw=keptline,line width=0.6pt] (21.0,-4) circle (0.17);
  \filldraw[fill=white,draw=keptline,line width=0.6pt] (22.0,-4) circle (0.17);
  \filldraw[fill=white,draw=keptline,line width=0.6pt] (23.0,-4) circle (0.17);
  \fill[black] (24.0,-4) circle (0.17);
  \draw[draw=rectA,line width=1.3pt,line join=round] (23.0,-1) -- (24.0,-2) -- (22.0,-4) -- (21.0,-3) -- cycle;
  \draw[rectA,line width=1.1pt] (21.0,-3) circle (0.3);
  \draw[rectA,line width=1.1pt] (22.0,-4) circle (0.3);
  \draw[draw=rectB,line width=1.3pt,line join=round,dash pattern=on 3pt off 2pt] (22.0,-1) -- (24.0,-2) -- (23.0,-4) -- (21.0,-3) -- cycle;
  \draw[rectB,line width=1.1pt] (21.0,-3) circle (0.38);
  \draw[rectB,line width=1.1pt] (23.0,-4) circle (0.38);
  \node[capt,anchor=north] at (21.0,-4.7) {(c) $\kA_4$ and $\kB_4$};
\end{tikzpicture}

\caption{Five of the weak tilted rectangles of the construction for $m=7$, with their shown corners ringed. (a) $D(1,1)$. (b) $\kA_1$ (solid) and $\kB_1$ (dashed). (c) $\kA_4$ (solid) and $\kB_4$ (dashed). No other rectangle shows exactly these points.}
\label{fig:fourtilted}
\end{figure}

\subsection{Widths four and five}

On the $4\times4$ and $4\times5$ grids the row bound says $m+4$, but the answer is $m+3$.

\begin{lemma}[equality in the row bound on four rows]\label{lem:fourequality}
Let $m\ge4$. A blackout of the $4\times m$ grid with $m+4$ points is axis-valid if and only if every column is nonempty, exactly four columns have two black points, their row pairs are the four edges of a four-cycle on the rows, and every other column has one black point. The number of such blackouts is
\begin{equation}\label{eq:fourcand}
3\,(m)_4\,4^{m-4},
\end{equation}
which is $72$, $1{,}440$, $17{,}280$, and $161{,}280$ for $m=4$, $5$, $6$, $7$. In particular every valid blackout with $m+4$ points has this form.
\end{lemma}

\begin{proof}
This is Theorem~\ref{thm:equality} with $h=4$ and $q_4=4$. A complete bipartite graph on four rows with parts of size two is a four-cycle, and there are three. To count, choose the cycle, a column for each of its four labeled edges in $(m)_4$ ways, and a black row in each other column.
\end{proof}

\begin{corollary}[widths four and five, computer-assisted]\label{cor:fourboundary}
We have $f(4,4)=7$ and $f(4,5)=8$.
\end{corollary}

\begin{proof}
The lower bounds are Lemma~\ref{lem:cross} with $h=4$. For the upper bounds, Lemma~\ref{lem:fourequality} lists the $72$ and $1{,}440$ blackouts with $m+4$ points that could be valid. The script \texttt{verify\_narrow.py} generates them from the lemma, generates every rectangle from pairs of equal diagonals with a common midpoint, and finds two rectangles with the same presentation in every candidate. The brute-force check of Section~\ref{sec:verify} tests every blackout of these two grids from the definition and confirms both values without the lemma.
\end{proof}

\subsection{No hidden rectangle}

Section~\ref{sec:conv} allows one hidden rectangle. On four rows of width at least four a valid blackout hides none. The proof writes each tilted rectangle as the symmetric difference of two other rectangles in the same four columns and applies Lemma~\ref{lem:diff}. For $1\le i\le m-3$ let $Q_i=\{(i,2),(i,3),(i+3,2),(i+3,3)\}$. Checking corners gives
\begin{equation}\label{eq:fourxor}
\kC_i=\kA_i\,\triangle\,Q_i,\qquad \kE_i=\kB_i\,\triangle\,Q_i,\qquad
D(i+1,1)=D(i,2)\,\triangle\,\kA_i,\qquad D(i+1,2)=D(i,1)\,\triangle\,\kE_i,
\end{equation}
where each rectangle stands for its set of corners. Since $X=Y\triangle Z$ is the same as $Y=X\triangle Z$, each identity can be solved for each of its three rectangles.

\begin{theorem}[no hidden rectangle]\label{thm:hidden}
Let $m\ge4$. A valid blackout of the $4\times m$ grid hides no rectangle, that is, every rectangle shows at least one corner. For $m=3$ this fails. The blackout $\{(1,2),(2,1),(3,2),(2,3)\}$ of the $4\times3$ grid is valid and hides the diamond $D(1,1)$.
\end{theorem}

\begin{proof}
Let $S$ be valid. By Lemma~\ref{lem:diff}, two rectangles have the same presentation exactly when the symmetric difference of their corner sets lies in $S$. So a hidden rectangle whose corner set is $C(R_1)\triangle C(R_2)$ for distinct rectangles $R_1,R_2$ contradicts validity, and we find such $R_1,R_2$ for every rectangle. If the axis-aligned rectangle on rows $r,s$ and columns $x,y$ is hidden, take a third column $w$. The rectangles on rows $r,s$ and columns $\{x,w\}$ and $\{y,w\}$ both show exactly the white points among $(w,r)$ and $(w,s)$. If $\kC_i$ is hidden, then $\kA_i\triangle Q_i\subseteq S$ by \eqref{eq:fourxor}, so $\kA_i$ and $Q_i$ have the same presentation. The identities $\kE_i=\kB_i\triangle Q_i$, $\kA_i=\kC_i\triangle Q_i$, and $\kB_i=\kE_i\triangle Q_i$ handle $\kE_i$, $\kA_i$, and $\kB_i$. A diamond occupies three columns, so for $m\ge4$ it has a column to its left or to its right. If $D(j,1)$ is hidden we use $D(j-1,2)\triangle\kA_{j-1}=D(j,1)$ when $j\ge2$ and $D(j+1,2)\triangle\kE_j=D(j,1)$ when $j\le m-3$. If $D(j,2)$ is hidden we use $D(j-1,1)\triangle\kE_{j-1}=D(j,2)$ when $j\ge2$ and $D(j+1,1)\triangle\kA_j=D(j,2)$ when $j\le m-3$. So no rectangle is hidden.

For $m=3$ the blackout $S$ is the corner set of $D(1,1)$, so $D(1,1)$ is hidden. The $4\times3$ grid has $18$ axis-aligned rectangles and the diamonds $D(1,1)$, $D(1,2)$. Ten of them show at least three corners and are safe by Lemma~\ref{lem:threecorners}. The other nine show the distinct pairs $\{(1,1),(3,1)\}$, $\{(1,3),(3,3)\}$, $\{(1,4),(3,4)\}$, $\{(1,1),(1,3)\}$, $\{(3,1),(3,3)\}$, $\{(1,1),(2,2)\}$, $\{(1,3),(2,2)\}$, $\{(2,2),(3,1)\}$, $\{(2,2),(3,3)\}$. So $S$ is valid.
\end{proof}

\subsection{Forbidden blocks}

In this subsection $m\ge6$, so by Theorem~\ref{thm:fourvalue} and Lemma~\ref{lem:fourequality} the optima are the valid blackouts with four double columns and singletons elsewhere. The pair set $F$ of such a blackout consists of the four row pairs of its double columns, the edges of a four-cycle on the rows. We name $F$ by the two parts of the cycle, $\{1,2\}\mid\{3,4\}$, $\{1,3\}\mid\{2,4\}$, or $\{1,4\}\mid\{2,3\}$, so that $F$ consists of the four pairs with one row in each part. The construction \eqref{eq:fourS} has the parts $\{1,4\}\mid\{2,3\}$.

A \emph{block} is a set of four consecutive singleton columns $x,\ldots,x+3$, and its \emph{word} is $w=w_1w_2w_3w_4$, where $w_k$ is the black row of column $x+k-1$. The tilted rectangles inside a block are the diamonds $D(x,1)$, $D(x+1,1)$, $D(x,2)$, $D(x+1,2)$ and the long rectangles $\kA_x$, $\kB_x$, $\kC_x$, $\kE_x$. Reading the corner rows of $\kC_x$, $\kA_x$, $\kE_x$, $\kB_x$ column by column gives the four \emph{template words} $2413$, $3412$, $2143$, $3142$. So $\kA_x$ is hidden in a block exactly when the word of the block is $3412$, and likewise for the other three. Three families of words matter.
\begin{enumerate}[(a)]
\item A \emph{swap word} is a template word with two of its letters exchanged.
\item An \emph{edit word for $F$} is a template word with one letter $r$ replaced by a row $s$ with $\{r,s\}\in F$.
\item A \emph{diamond word for $F$} is a word with $w_1=w_3=r$ or $w_2=w_4=r$, where $r\in\{2,3\}$ and $\{r-1,r+1\}\in F$.
\end{enumerate}
An axis-aligned rectangle with the same presentation as a tilted rectangle is an \emph{axis rival} of it. The word $w$ of a block of a blackout of the equality form with pair set $F$ is \emph{admissible for $F$} if no tilted rectangle inside the block is hidden or has an axis rival. By Lemma~\ref{lem:rival} this depends only on $w$ and $F$, not on where the double columns are.

\begin{lemma}[forbidden blocks]\label{lem:blocks}
Let $m\ge6$ and let $S$ be a valid blackout of the $4\times m$ grid with $m+4$ points and pair set $F$. Then the word of every block of $S$, if it has one, is admissible for $F$. Moreover, with the three cases ordered as $\{1,2\}\mid\{3,4\}$, $\{1,3\}\mid\{2,4\}$, $\{1,4\}\mid\{2,3\}$,
\begin{enumerate}[(a)]
\item the swap words are the $20$ permutations of $1234$ other than $1234$, $1324$, $4231$, and $4321$,
\item there are $24$, $24$, or $32$ edit words for $F$, and no edit word is a swap word,
\item there are $60$, $0$, or $60$ diamond words for $F$, of which $4$, $0$, or $8$ are also edit words, and no diamond word is a swap word,
\item a word is admissible for $F$ if and only if it belongs to none of the three families, so the number of admissible words is $156$, $212$, or $152$ (Table~\ref{tab:fourblocks}).
\end{enumerate}
\end{lemma}

\begin{table}[t]
\centering\small
\begin{tabular}{@{}lrrrrrr@{}}
\toprule
parts of $F$ & swap & edit & diamond & overlap & excluded & admissible\\
\midrule
$\{1,2\}\mid\{3,4\}$ & $20$ & $24$ & $60$ & $4$ & $100$ & $156$\\
$\{1,3\}\mid\{2,4\}$ & $20$ & $24$ & $0$ & $0$ & $44$ & $212$\\
$\{1,4\}\mid\{2,3\}$ & $20$ & $32$ & $60$ & $8$ & $104$ & $152$\\
\bottomrule
\end{tabular}
\caption{The words of a block of four singleton columns, out of $4^4=256$, by Lemma~\ref{lem:blocks}. The overlap column counts the words that are both edit words and diamond words. The swap words are disjoint from the other two families.}
\label{tab:fourblocks}
\end{table}

\begin{proof}
A valid blackout has no hidden rectangle (Theorem~\ref{thm:hidden}) and no axis rival, so the word of every block is admissible. Throughout, $T$ is one of $\kC_x$, $\kA_x$, $\kE_x$, $\kB_x$ with template word $t=t_1t_2t_3t_4$, so the corner of $T$ in column $x+k-1$ lies in row $t_k$.

\emph{Step 1. Swap words.} Let $w$ be $t$ with the letters in positions $p\ne q$ exchanged. Then $T$ has black corners in the other two columns and white corners $(x+p-1,t_p)$ and $(x+q-1,t_q)$, since these columns are black in rows $t_q$ and $t_p$. The axis-aligned rectangle on the two white corners has the other corners $(x+p-1,t_q)$ and $(x+q-1,t_p)$, which are black, so it is an axis rival. For the count, the template words are exactly the permutations with both end letters in $\{2,3\}$. Exchanging two end letters or two middle letters gives a template word again, and exchanging an end letter with a middle letter gives a permutation with exactly one end letter in $\{2,3\}$, and every such permutation arises this way, from the template word obtained by exchanging its misplaced end letter with its misplaced middle letter. So the swap words are the permutations with at least one end letter in $\{2,3\}$, all but $1234$, $1324$, $4231$, $4321$. This proves (a).

\emph{Step 2. Edit words.} Let $w$ be $t$ with $r=t_p$ replaced by $s$, where $\{r,s\}\in F$. Then $T$ shows only $(x+p-1,r)$, and column $x+p-1$ is black in row $s$. Let $d$ be the double column with black rows $r,s$. The axis-aligned rectangle on columns $x+p-1,d$ and rows $r,s$ also shows only $(x+p-1,r)$, so it is an axis rival, however far away $d$ is. For the count, each template word has $4\cdot2=8$ edit words, so there are $32$ pairs of a template word and an edit word of it. A word arises from two template words $t\ne t'$ only when they differ in exactly two positions, carrying two letters $r,s$ in opposite orders, and the word has $r$ in both positions or $s$ in both. These two words are edit words of both exactly when $\{r,s\}\in F$. The pairs at distance two are $2413,3412$ and $2143,3142$ with $\{r,s\}=\{2,3\}$, and $2413,2143$ and $3412,3142$ with $\{r,s\}=\{1,4\}$. So there are $32-4[\{2,3\}\in F]-4[\{1,4\}\in F]$ edit words, which is $24$, $24$, or $32$. An edit word repeats a letter, so it is not a swap word. This proves (b).

\emph{Step 3. Diamond words.} Let $w_{j+1}=w_{j+3}=r$ with $j\in\{0,1\}$, $r\in\{2,3\}$, and $\{r-1,r+1\}\in F$. The diamond $D(x+j,r-1)$ has the black tips $(x+j,r)$, $(x+j+2,r)$ and the middle corners $(c,r-1)$, $(c,r+1)$ in the singleton column $c=x+j+1$. Let $d$ be the double column with black rows $r-1,r+1$. The diamond and the axis-aligned rectangle on rows $r-1,r+1$ and columns $c,d$ both show exactly the white points among the two middle corners, so the latter is an axis rival. For the count, in the first and third cases both $\{1,3\}$ and $\{2,4\}$ lie in $F$, so $r$ may be $2$ or $3$. For each $j$ there are $2\cdot4^2=32$ words, and $4$ words have $w_1=w_3$ and $w_2=w_4$ with both values in $\{2,3\}$, so there are $60$ diamond words. In the middle case there are none. A diamond word repeats a letter, so it is not a swap word. An edit word with $w_1=w_3$ is a template word $t$ with its first letter copied to position $3$, since copying $t_3\in\{1,4\}$ to position $1$ gives $r\notin\{2,3\}$. It is an edit word when $\{t_1,t_3\}\in F$ and a diamond word when $\{t_1-1,t_1+1\}\in F$. Likewise an edit word with $w_2=w_4$ is $t$ with its fourth letter copied to position $2$, with the conditions $\{t_2,t_4\}\in F$ and $\{t_4-1,t_4+1\}\in F$. The eight candidates $2423$, $2313$, $3432$, $3212$, $2123$, $2343$, $3132$, $3242$ are distinct. For $\{1,2\}\mid\{3,4\}$ the words $3432$, $3212$, $2123$, $2343$ satisfy both conditions, for $\{1,4\}\mid\{2,3\}$ all eight do, and for $\{1,3\}\mid\{2,4\}$ there are no diamond words. This proves (c).

\emph{Step 4. Nothing else is excluded.} Let some tilted rectangle inside the block be hidden or have an axis rival. If it is a long rectangle $T$, its corners lie in distinct rows and columns. Four black corners mean $w=t$. With three, $T$ shows one point $(x+p-1,t_p)$, and by Lemma~\ref{lem:rival} a rival exists exactly when the column is black in a row $s$ with $\{t_p,s\}\in F$, so $w$ is an edit word. With two, $T$ shows an oblique pair, and by Lemma~\ref{lem:rival} a rival exists exactly when each column is black in the row of the other point, so $w$ is a swap word. With at most one, $T$ is safe by Lemma~\ref{lem:threecorners}. If it is a diamond with tips in row $r$ and middle column $c$, at most one middle corner is black, since $c$ is a singleton. If both tips are black, the diamond shows one or both middle corners, in column $c$, and by Lemma~\ref{lem:rival} a rival exists exactly when $\{r-1,r+1\}\in F$, so $w$ is a diamond word. If one tip and one middle corner are black, a rival of the shown oblique pair would need $(c,r)$ black, which is false. Otherwise the diamond shows at least three corners. So $w$ belongs to one of the three families, and the number of admissible words is $256-20-|E|-|D|+|E\cap D|$ with $E$, $D$ the edit and diamond families, which gives $156$, $212$, $152$. This proves (d).
\end{proof}

\begin{theorem}[forbidden-block bound]\label{thm:foursurvival}
For every $m\ge6$,
\begin{equation}\label{eq:survival}
\begin{split}
a_4(m)&\le(m)_4\,4^{m-4}\Bigl[\bigl(\tfrac{39}{64}\bigr)^K+\bigl(\tfrac{53}{64}\bigr)^K+\bigl(\tfrac{19}{32}\bigr)^K\Bigr]\\
&\le3\,(m)_4\,4^{m-4}\bigl(\tfrac{53}{64}\bigr)^K ,
\end{split}
\end{equation}
where $K=\max(0,\lfloor m/4\rfloor-4)$. Since $4^4\cdot\frac{53}{64}=212$, this gives $a_4(m)=O(m^4\,212^{m/4})$.
\end{theorem}

For $6\le m\le19$ we have $K=0$, so the bound says something new only from $m=20$ on. Its point is the base $212^{1/4}=3.8157\ldots<4$. The fraction of the axis-valid candidates that are valid tends to zero exponentially, while on three rows it tends to $4/27$ (Corollary~\ref{cor:threeasymp}). Corollary~\ref{cor:fourgrowth} computes the true growth rate.

\begin{proof}
Fix the parts of $F$, the four double columns, and the pair carried by each, in one of $(m)_4$ ways for each $F$. Cut the columns into the $\lfloor m/4\rfloor$ disjoint groups $\{4k-3,\ldots,4k\}$ and a remainder. At most four groups contain a double column, so at least $K$ groups are blocks. By Lemma~\ref{lem:blocks} a block has at most $c_F$ words, where $c_F$ is $156$, $212$, or $152$, and every other singleton column has at most $4$ values. So if $N'\ge K$ groups are blocks, this support has at most $4^{m-4}(c_F/256)^{N'}\le4^{m-4}(c_F/256)^K$ valid blackouts. Summing over the supports and the pair sets, with $156/256=39/64$, $212/256=53/64$, and $152/256=19/32$, gives the first inequality of \eqref{eq:survival}, and the second follows from $39/64<53/64$ and $19/32<53/64$.
\end{proof}

\subsection{Six columns of memory}

Theorem~\ref{thm:foursurvival} uses one necessary condition. To count exactly we need a test that is also sufficient and reads the columns from left to right with bounded memory. Four columns are not enough.

\begin{example}\label{ex:fourwidtheight}
Take the blackout of the $4\times8$ grid whose columns are black in the rows $\{1,2\}$, $\{2,3\}$, $\{3\}$, $\{4\}$, $\{3,4\}$, $\{1\}$, $\{3\}$, $\{1,4\}$, as in Figure~\ref{fig:widtheight}. It has $12=m+4$ points, its double columns carry the pairs of $\{1,3\}\mid\{2,4\}$, and it is axis-valid by Lemma~\ref{lem:fourequality}. The diamond $T=D(3,2)=\{(3,3),(4,2),(5,3),(4,4)\}$ and the long rectangle $U=\kC_4=\{(4,2),(5,4),(6,1),(7,3)\}$ both show only $(4,2)$, so the blackout is not valid. Column $4$ is black in row $4$ and $\{2,4\}\notin F$, so this presentation has no axis rival (Lemma~\ref{lem:rival}), and among the $200$ rectangles of the grid $T,U$ is the only pair with the same presentation. Their union occupies columns $3$ to $7$. So a test that looks at one tilted rectangle at a time, or at pairs of tilted rectangles inside four consecutive columns, accepts this blackout.
\end{example}

\begin{figure}[t]
\centering
\begin{tikzpicture}[x=0.500cm,y=0.500cm,baseline=(current bounding box.north)]
  \draw[gridcol] (1,-1) grid (8,-4);
  \fill[black] (1,-1) circle (0.17);
  \filldraw[fill=white,draw=keptline,line width=0.6pt] (2,-1) circle (0.17);
  \filldraw[fill=white,draw=keptline,line width=0.6pt] (3,-1) circle (0.17);
  \filldraw[fill=white,draw=keptline,line width=0.6pt] (4,-1) circle (0.17);
  \filldraw[fill=white,draw=keptline,line width=0.6pt] (5,-1) circle (0.17);
  \fill[black] (6,-1) circle (0.17);
  \filldraw[fill=white,draw=keptline,line width=0.6pt] (7,-1) circle (0.17);
  \fill[black] (8,-1) circle (0.17);
  \fill[black] (1,-2) circle (0.17);
  \fill[black] (2,-2) circle (0.17);
  \filldraw[fill=white,draw=keptline,line width=0.6pt] (3,-2) circle (0.17);
  \filldraw[fill=white,draw=keptline,line width=0.6pt] (4,-2) circle (0.17);
  \filldraw[fill=white,draw=keptline,line width=0.6pt] (5,-2) circle (0.17);
  \filldraw[fill=white,draw=keptline,line width=0.6pt] (6,-2) circle (0.17);
  \filldraw[fill=white,draw=keptline,line width=0.6pt] (7,-2) circle (0.17);
  \filldraw[fill=white,draw=keptline,line width=0.6pt] (8,-2) circle (0.17);
  \filldraw[fill=white,draw=keptline,line width=0.6pt] (1,-3) circle (0.17);
  \fill[black] (2,-3) circle (0.17);
  \fill[black] (3,-3) circle (0.17);
  \filldraw[fill=white,draw=keptline,line width=0.6pt] (4,-3) circle (0.17);
  \fill[black] (5,-3) circle (0.17);
  \filldraw[fill=white,draw=keptline,line width=0.6pt] (6,-3) circle (0.17);
  \fill[black] (7,-3) circle (0.17);
  \filldraw[fill=white,draw=keptline,line width=0.6pt] (8,-3) circle (0.17);
  \filldraw[fill=white,draw=keptline,line width=0.6pt] (1,-4) circle (0.17);
  \filldraw[fill=white,draw=keptline,line width=0.6pt] (2,-4) circle (0.17);
  \filldraw[fill=white,draw=keptline,line width=0.6pt] (3,-4) circle (0.17);
  \fill[black] (4,-4) circle (0.17);
  \fill[black] (5,-4) circle (0.17);
  \filldraw[fill=white,draw=keptline,line width=0.6pt] (6,-4) circle (0.17);
  \filldraw[fill=white,draw=keptline,line width=0.6pt] (7,-4) circle (0.17);
  \fill[black] (8,-4) circle (0.17);
  \node[lab] at (1,-0.35) {1};
  \node[lab] at (2,-0.35) {2};
  \node[lab] at (3,-0.35) {3};
  \node[lab] at (4,-0.35) {4};
  \node[lab] at (5,-0.35) {5};
  \node[lab] at (6,-0.35) {6};
  \node[lab] at (7,-0.35) {7};
  \node[lab] at (8,-0.35) {8};
  \node[lab] at (0.35,-1) {1};
  \node[lab] at (0.35,-2) {2};
  \node[lab] at (0.35,-3) {3};
  \node[lab] at (0.35,-4) {4};
  \draw[draw=rectA,line width=1.3pt,line join=round] (3,-3) -- (4,-2) -- (5,-3) -- (4,-4) -- cycle;
  \draw[draw=rectB,line width=1.3pt,line join=round,dash pattern=on 3pt off 2pt] (4,-2) -- (5,-4) -- (7,-3) -- (6,-1) -- cycle;
  \draw[rectA,line width=1.1pt] (4,-2) circle (0.3);
  \draw[rectB,line width=1.1pt] (4,-2) circle (0.38);
  \node[tiny,anchor=south west,text=rectA] at (2.55,-3.45) {$T$};
  \node[tiny,anchor=north east,text=rectB] at (7.45,-2.55) {$U$};
\end{tikzpicture}

\caption{The blackout of the $4\times8$ grid with black rows $\{1,2\}$, $\{2,3\}$, $\{3\}$, $\{4\}$, $\{3,4\}$, $\{1\}$, $\{3\}$, $\{1,4\}$ in columns $1$ to $8$. The diamond $T=D(3,2)$ (solid) and the long rectangle $U=\kC_4$ (dashed) each show only the ringed point $(4,2)$.}
\label{fig:widtheight}
\end{figure}

\begin{theorem}[six columns of memory]\label{thm:fourmemory}
Fix one of the three pair sets $F$ and let $m\ge4$. The letter of a column is its set of black rows. Read a blackout of the $4\times m$ grid column by column as a word over the eight letters $\{1\}$, $\{2\}$, $\{3\}$, $\{4\}$, and the four pairs in $F$. Reject at once a column with another letter, and a pair of $F$ that occurred before. After each column, consider the last at most seven columns read, and reject if inside them
\begin{enumerate}[(a)]
\item some tilted rectangle is hidden,
\item some tilted rectangle has an axis rival in the finished board, which by Lemma~\ref{lem:rival} is decided by these columns and $F$, or
\item two distinct tilted rectangles have the same nonempty presentation.
\end{enumerate}
Accept the word at the end if it was never rejected and every pair of $F$ occurred. Then the accepted words are exactly the valid blackouts with $m+4$ points and pair set $F$. The test after each column uses only the six previous columns, the new column, the pairs of $F$ used so far, and $F$ itself.
\end{theorem}

We do not claim that six columns are the least possible memory.

\begin{proof}
An accepted word has four double columns carrying the four pairs of $F$ and singletons elsewhere, so it has $m+4$ points and is axis-valid by Lemma~\ref{lem:fourequality}. Rule (a) excludes hidden tilted rectangles, and a hidden axis-aligned rectangle would be a four-cycle of $G_S$, so no rectangle is hidden. Suppose two distinct rectangles have the same presentation. They are not both axis-aligned. If one is axis-aligned, the other is a tilted rectangle with an axis rival, which lies inside four consecutive columns, so rule (b) rejects when its last column is read. If both are tilted, they share a shown corner and with it a column $c$, so both lie inside the columns $c-3$ to $c+3$, and rule (c) rejects when the later of their last columns is read. So accepted words are valid. Conversely let $S$ be valid with $m+4$ points and pair set $F$. By Lemma~\ref{lem:fourequality} its word is over the eight letters, no pair repeats, and every pair occurs. Rule (a) never fires by Theorem~\ref{thm:hidden}, and rules (b) and (c) fire only when a tilted rectangle shares its presentation with another rectangle of the finished board, which does not happen. So $S$ is accepted. After a column is read, the window is that column and the six before it, which gives the memory.
\end{proof}

The test is a finite automaton, so the counting is a transfer-matrix computation (Flajolet and Sedgewick~\cite[Section~V.5]{flajolet2009}), and the four double columns cut every word into five runs of singletons.

\begin{theorem}[five-run formula]\label{thm:fourgf}
Fix a pair set $F$. The state of the automaton of Theorem~\ref{thm:fourmemory} after a word is read is the last six letters of the word, or the whole word when it is shorter. Let $\Sigma_F$ be the set of states reached from the empty word without a rejection by the rules (a) to (c), the rule against a repeated pair being left to the ordering below. Let $S_F$ be the matrix on $\Sigma_F$ whose $(u,v)$ entry is the number of singleton letters that lead from $u$ to $v$ without rejection, and for $e\in F$ let $D_{F,e}$ be the same matrix for the single letter $e$. Let $\alpha_F$ select the empty word, let $\mathbf1$ be the all-ones column vector, and put $R_F(z)=(I-zS_F)^{-1}$. Then
\begin{equation}\label{eq:fiverun}
\begin{split}
\sum_{m\ge6}a_4(m)\,z^m=z^4\sum_F\ \sum_{(e_1,e_2,e_3,e_4)}&\alpha_F\,R_F(z)\,D_{F,e_1}\,R_F(z)\,D_{F,e_2}\\
&\cdot R_F(z)\,D_{F,e_3}\,R_F(z)\,D_{F,e_4}\,R_F(z)\,\mathbf1,
\end{split}
\end{equation}
where the inner sum runs over the $24$ orderings of the four pairs of $F$. In particular $\sum_{m\ge6}a_4(m)z^m$ is a rational function of $z$.
\end{theorem}

The right side of \eqref{eq:fiverun} has the coefficient $0$ at $z^4$ and $z^5$ (Corollary~\ref{cor:fourboundary}), while $a_4(4)=316$ and $a_4(5)=1{,}810$ count blackouts with $m+3$ points.

\begin{proof}
By Theorems~\ref{thm:fourvalue} and~\ref{thm:fourmemory}, the words of length $m\ge6$ accepted for the three pair sets are the $a_4(m)$ optima. Every accepted word is $s_0e_1s_1e_2s_2e_3s_3e_4s_4$ for exactly one ordering $(e_1,\ldots,e_4)$ of $F$ and one choice of five possibly empty runs $s_0,\ldots,s_4$ of singleton letters. The automaton is deterministic, the entry $(u,v)$ of $S_F^k$ counts the runs of $k$ singleton letters that lead from $u$ to $v$ without rejection, and $D_{F,e}$ records the step on $e$. So the accepted words with the ordering $(e_1,\ldots,e_4)$ and run lengths $k_0,\ldots,k_4$ number $\alpha_FS_F^{k_0}D_{F,e_1}S_F^{k_1}\cdots D_{F,e_4}S_F^{k_4}\mathbf1$, where $\mathbf1$ counts every surviving path because the condition that every pair occurred is built into the ordering. Summing over the run lengths with the weight $z^{k_0+\cdots+k_4+4}$ and using $\sum_{k\ge0}z^kS_F^k=R_F(z)$ gives \eqref{eq:fiverun}. The entries of $R_F(z)=\operatorname{adj}(I-zS_F)/\det(I-zS_F)$ are rational, so the sum is rational.
\end{proof}

To compute the values we do not invert $I-zS_F$. For $U\subseteq F$ let $v_U(n)$ be the row vector indexed by $\Sigma_F$ whose entry at $u$ counts the words of length $n$ that end in the state $u$, were never rejected, and use exactly the pairs in $U$. Then $v_\emptyset(0)=\alpha_F$, the other $v_U(0)$ vanish, and
\begin{equation}\label{eq:fourrec}
v_U(n+1)=v_U(n)\,S_F+\sum_{e\in U}v_{U\setminus\{e\}}(n)\,D_{F,e},\qquad a_4(n)=\sum_F v_F(n)\,\mathbf1\quad(n\ge6),
\end{equation}
according to whether the last letter is a singleton or a new pair.

\subsection{Row-bound attainment at every height}

The $20$ swap words of Lemma~\ref{lem:blocks}(a) are forbidden by the block alone, without the pair set or the height. They give a bound at every height.

\begin{corollary}[rarity of row-bound attainment]\label{cor:rarity}
Let $h\ge4$ and $q_h=\lfloor h^2/4\rfloor$. For $m\ge3$ let $b_h(m)$ be the number of valid blackouts of the $h\times m$ grid with $m+q_h$ points, and let
\[
C_h(m)=p_h\,(m)_{q_h}\,h^{m-q_h},\qquad
p_h=\begin{cases}\tfrac12\binom{h}{h/2}&\text{if $h$ is even,}\\ \binom{h}{\lfloor h/2\rfloor}&\text{if $h$ is odd,}\end{cases}
\]
be the number of axis-valid blackouts with $m+q_h$ points. Then $b_h(m)=0$ for $3\le m<q_h$, and for $m\ge q_h$
\begin{equation}\label{eq:rarity}
\frac{b_h(m)}{C_h(m)}\le\Bigl(1-\frac{20(h-3)}{h^4}\Bigr)^{\max(0,\lfloor m/4\rfloor-q_h)} .
\end{equation}
In particular $b_h(m)=O_h\bigl(m^{q_h}(h^4-20h+60)^{m/4}\bigr)$. For $h=5$, since $f(5,m)=m+6$ for $m\ge21$ (Theorem~\ref{thm:fiverow}), this gives
\[
\frac{a_5(m)}{10\,(m)_6\,5^{m-6}}\le\Bigl(\frac{117}{125}\Bigr)^{\max(0,\lfloor m/4\rfloor-6)}\quad(m\ge21),\qquad a_5(m)=O(m^6\,585^{m/4}).
\]
\end{corollary}

At $h=4$ the corollary gives the factor $59/64$ per block, and Theorem~\ref{thm:foursurvival} improves it to $53/64$ by using $F$. The factor $20(h-3)$ is not optimized.

\begin{proof}
By Theorem~\ref{thm:equality} the axis-valid blackouts with $m+q_h$ points are those with $q_h$ double columns whose row pairs are the edges of a balanced complete bipartite graph on the rows, each edge once, and singletons elsewhere, $C_h(m)$ of them for $m\ge q_h$ and none below. A valid blackout is axis-valid, so $b_h(m)\le C_h(m)$, and $b_h(m)=0$ for $3\le m<q_h$ (Corollary~\ref{cor:small-m}). For $0\le j\le h-4$ add $j$ to every letter of the $20$ swap words, which gives $20$ words over the rows $j+1,\ldots,j+4$. If a block of four consecutive singleton columns has one of them as its word, Step 1 of the proof of Lemma~\ref{lem:blocks} applies on these four rows. So a long tilted rectangle shows an oblique pair whose axis-aligned rectangle has its other two corners black, and this uses only the four columns of the block and no pair set. Words for different $j$ differ, since each uses all four of its rows. So at least $20(h-3)$ of the $h^4$ words of a block are forbidden in every valid blackout. Now fix one of the $p_h(m)_{q_h}$ choices of the double columns and their pairs. Of the $\lfloor m/4\rfloor$ disjoint groups of columns $4k-3,\ldots,4k$ at most $q_h$ contain a double column. So at least $\max(0,\lfloor m/4\rfloor-q_h)$ groups are blocks, each with at most $h^4-20(h-3)$ words, and every other singleton column has at most $h$ values. Multiplying and summing over the supports gives \eqref{eq:rarity}, and $h^4-20(h-3)=h^4-20h+60$ gives the growth bound. For five rows, $p_5=10$, $q_5=6$, $5^4-40=585$, $1-40/625=117/125$, and $a_5(m)=b_5(m)$ for $m\ge21$ by Theorem~\ref{thm:fiverow}.
\end{proof}

\subsection{The first values and the generating function}

\begin{proposition}[small widths, computer-assisted]\label{prop:fourvalues}
The values $a_4(m)$ for $4\le m\le7$ in the row $h=4$ of \seqnum{A399580} are $a_4(4)=316$, $a_4(5)=1{,}810$, $a_4(6)=26$, and $a_4(7)=292$. Order the pair sets as $\{1,2\}\mid\{3,4\}$, $\{1,3\}\mid\{2,4\}$, $\{1,4\}\mid\{2,3\}$. Then the $26$ optima of width $6$ split as $0$, $20$, $6$, and the $292$ optima of width $7$ split as $4$, $256$, $32$.
\end{proposition}

\begin{proof}
By Corollary~\ref{cor:fourboundary} the optima of widths $4$ and $5$ have $7$ and $8$ points. The ILP program and the brute-force check of Section~\ref{sec:verify} both find $316$ and $1{,}810$ of them, and both find $26$ and $292$ at widths $6$ and $7$. By Lemma~\ref{lem:fourequality} the optima of widths $6$ and $7$ are among the $17{,}280$ and $161{,}280$ candidates of \eqref{eq:fourcand}, and \texttt{four\_row\_direct.py} tests each candidate against the $110$ and $152$ rectangles of the grid and leaves $26$ and $292$, with the stated splits. The recurrence \eqref{eq:fourrec} on the machine of \texttt{four\_row\_counts.py} reproduces both values and both splits. That machine is another implementation of Theorem~\ref{thm:fourmemory}, storing the matched prefixes of a finite list of forbidden patterns instead of six columns, with $383$, $487$, and $307$ reachable states for the three pair sets, computed sizes rather than minimal ones. Section~\ref{sec:verify} lists the scripts.
\end{proof}

The recurrence \eqref{eq:fourrec} on that machine gives $a_4(m)$ for every $m\le1500$ in a few seconds, and from these values the rational function of Theorem~\ref{thm:fourgf} can be found and certified.

\begin{theorem}[the four-row generating function, computer-assisted]\label{thm:fourgfexact}
There are coprime integer polynomials $P_4$ of degree $444$ and $Q_4$ of degree $420$ with $Q_4(0)=1$ such that
\[
\sum_{m\ge6}a_4(m)\,z^m=\frac{P_4(z)}{Q_4(z)},\qquad P_4(z)=26z^6-826z^7+12052z^8-107652z^9+\cdots .
\]
The denominator factors over the integers as
\[
Q_4=(1-z)^2(1+z)^2(1+z^2-z^5)^3(1-z^2-z^5-2z^6)^3(1+3z^2-2z^3+2z^4-z^5)^3\,g_{11}^5\,g_{17}^3\,g_{19}^3\,g_{20}^5\,g_{21}^5,
\]
where $g_d$ is an irreducible polynomial of degree $d$ with constant term $1$, and $g_{20}(z)=z^{20}p(1/z)$ for the polynomial $p$ of Corollary~\ref{cor:fourgrowth}. With $Q_4(z)=\sum_{j=0}^{420}q_jz^j$,
\begin{equation}\label{eq:fourrecur}
\sum_{j=0}^{420}q_j\,a_4(n-j)=0\qquad\text{for every }n\ge445,
\end{equation}
and \eqref{eq:fourrecur} fails at $n=444$. The polynomials and the factors $g_d$ are in the data file \texttt{four\_row\_recurrence.json}.
\end{theorem}

\begin{proof}
The three supports are treated separately and the three rational functions are added. For a pair set $F$ let $T_F$ be the transfer operator of the product of the machine of Proposition~\ref{prop:fourvalues} with the register of used pairs, with $16\cdot383=6{,}128$, $7{,}792$, and $4{,}912$ states, so that $v(n)=v(0)T_F^n$ is the vector of \eqref{eq:fourrec} and the count of width $n$ is the sum of its accepting coordinates. The script \texttt{four\_row\_recurrence.py} runs the Berlekamp--Massey algorithm on $a_4(6),\ldots,a_4(1500)$ modulo eight primes near $2^{62}$, lifts the result to the integers, and checks it exactly on all $1{,}495$ terms. This gives $Q_4$ and $P_4=Q_4\sum_{m\ge6}a_4(m)z^m$, of degree exactly $444$, and a gcd computation shows that they are coprime. The certificate is all-coordinate annihilation. The script checks $\sum_{j=0}^{420}q_j\,v(K-j)=0$ in every coordinate at $K=436$, $435$, and $445$ for the three supports. Applying $T_F^r$ gives the identity at every $K+r$, and summing the accepting coordinates gives \eqref{eq:fourrecur} for every $n\ge445$. The failure at $n=444$ is an exact integer check, forced because for coprime $P,Q$ with $\deg P\ge\deg Q$ the recurrence starts at $\deg P+1$. The factorization is by sympy, and \texttt{four\_row\_recurrence\_check.py} rechecks with its own integer arithmetic that the factors multiply to $Q_4$, that each is irreducible modulo a prime, and the threshold $445$. A second certificate, by Cayley--Hamilton windows on each product automaton, is \texttt{four\_row\_window\_check.py}. Section~\ref{sec:verify} lists the scripts.
\end{proof}

\begin{corollary}[growth, computer-assisted]\label{cor:fourgrowth}
Let $\rho=3.351996777310\ldots$ be the largest real root of the irreducible polynomial
\begin{multline*}
p(x)=x^{20}-4x^{19}+2x^{18}+x^{17}-3x^{16}+5x^{15}+x^{14}-3x^{13}+7x^{12}+2x^{11}+30x^{10}\\
-7x^{9}+10x^{8}+12x^{7}+2x^{6}+17x^{5}+26x^{4}+19x^{3}-26x^{2}-4x+8 .
\end{multline*}
Then $a_4(m)\sim c\,m^4\rho^m$ with $c=8.1611\times10^{-5}$. The support $\{1,3\}\mid\{2,4\}$ alone carries the factor $g_{20}$, so its optima grow like $c\,m^4\rho^m$, while those of $\{1,2\}\mid\{3,4\}$ and $\{1,4\}\mid\{2,3\}$ grow like constants times $m^4\,2.9320^m$ and $m^4\,2.8569^m$. The survival ratio $a_4(m)/(3(m)_44^{m-4})$ rises from $m=6$ to $m=7$ and then decreases strictly, checked exactly through $m=1500$.
\end{corollary}

\begin{proof}
The roots of $Q_4$ of smallest modulus $1/\rho=0.29833\ldots$ are the five copies of the positive real root of $g_{20}$, and every other root of $Q_4$ has modulus at least $0.3410$. Since $P_4$ and $Q_4$ are coprime, $P_4/Q_4$ has a pole of order exactly $5$ at $1/\rho$ and no other singularity of that modulus, so $a_4(m)\sim c\,m^4\rho^m$ with $c$ the leading coefficient of the partial fraction at the pole (Flajolet and Sedgewick~\cite[Section~IV.5]{flajolet2009}). The denominators of the three supports are $g_{19}^3g_{21}^5$, then $(1+z)^2g_{17}^3g_{20}^5$, and then the five small factors times $g_{11}^5$, each with a unique dominant root, real and positive, of multiplicity $5$, which gives the three rates. The root moduli and $c$ are computed to high precision by \texttt{four\_row\_rho\_check.py}. The survival ratio decreases from $m$ to $m+1$ exactly when $a_4(m+1)(m-3)<4a_4(m)(m+1)$, which \texttt{four\_row\_recurrence\_check.py} checks in exact integers for $7\le m<1500$.
\end{proof}

Before the computation the ledger recorded the guesses $\rho\ge3.4$ and a recurrence order between $20$ and $300$, and both were wrong. The forbidden-block bound of Theorem~\ref{thm:foursurvival} now reads as the proved upper bound $212^{1/4}=3.8158\ldots$ on the base, against the computed $3.3520$. Table~\ref{tab:fourvalues} lists $a_4(m)$ for $4\le m\le20$ with the counts by support. The values of $f(4,m)$ form the row $h=4$ of \seqnum{A397315}.

\begin{table}[t]
\centering\footnotesize
\begin{tabular}{@{}rrrrrr@{}}
\toprule
$m$ & $f(4,m)$ & $a_4(m)$ & $\{1,2\}\mid\{3,4\}$ & $\{1,3\}\mid\{2,4\}$ & $\{1,4\}\mid\{2,3\}$\\
\midrule
$4$ & $7$ & $316$ & -- & -- & --\\
$5$ & $8$ & $1{,}810$ & -- & -- & --\\
$6$ & $10$ & $26$ & $0$ & $20$ & $6$\\
$7$ & $11$ & $292$ & $4$ & $256$ & $32$\\
$8$ & $12$ & $2{,}118$ & $128$ & $1{,}780$ & $210$\\
$9$ & $13$ & $13{,}332$ & $920$ & $11{,}128$ & $1{,}284$\\
$10$ & $14$ & $75{,}398$ & $4{,}852$ & $63{,}540$ & $7{,}006$\\
\midrule
$11$ & $15$ & $402{,}656$ & $24{,}112$ & $346{,}320$ & $32{,}224$\\
$12$ & $16$ & $2{,}037{,}848$ & $114{,}424$ & $1{,}785{,}768$ & $137{,}656$\\
$13$ & $17$ & $9{,}894{,}016$ & $509{,}204$ & $8{,}818{,}256$ & $566{,}556$\\
$14$ & $18$ & $46{,}432{,}230$ & $2{,}152{,}908$ & $42{,}026{,}616$ & $2{,}252{,}706$\\
$15$ & $19$ & $211{,}990{,}020$ & $8{,}809{,}148$ & $194{,}503{,}368$ & $8{,}677{,}504$\\
$16$ & $20$ & $945{,}890{,}248$ & $35{,}133{,}016$ & $878{,}139{,}288$ & $32{,}617{,}944$\\
$17$ & $21$ & $4{,}138{,}161{,}480$ & $136{,}882{,}568$ & $3{,}881{,}050{,}516$ & $120{,}228{,}396$\\
$18$ & $22$ & $17{,}797{,}763{,}512$ & $522{,}435{,}792$ & $16{,}839{,}724{,}240$ & $435{,}603{,}480$\\
$19$ & $23$ & $75{,}413{,}531{,}108$ & $1{,}959{,}334{,}512$ & $71{,}899{,}462{,}068$ & $1{,}554{,}734{,}528$\\
$20$ & $24$ & $315{,}374{,}655{,}654$ & $7{,}238{,}543{,}584$ & $302{,}658{,}629{,}668$ & $5{,}477{,}482{,}402$\\
\bottomrule
\end{tabular}
\caption{The four-row values, with the optima of widths $6$ to $20$ counted by the parts of their pair set. The maxima are Corollary~\ref{cor:fourboundary} and Theorem~\ref{thm:fourvalue}, and the counts are computer-assisted. Widths $4$ to $7$ are the ILP program and the brute-force check of Section~\ref{sec:verify}. Widths $6$ to $10$ are the machine of Proposition~\ref{prop:fourvalues} and, independently, the definitional enumeration of every equality candidate against every rectangle by \texttt{four\_row\_direct.py}. Widths $11$ to $20$ are the machine alone, and \texttt{four\_row\_counts.json} continues the table to $m=1500$.}
\label{tab:fourvalues}
\end{table}
